\documentclass[11pt]{ctexart}
\usepackage[a4paper,margin=1in]{geometry}
\usepackage{graphicx}
\usepackage{xcolor}
\usepackage{hyperref}
\definecolor{refgray}{gray}{0.45}
\title{\textbf{市场最新观点汇总}\\[0.4em]{\large 市场主线：供给冲击与结构性分化主导资产定价，AI需求持续超预期但盈利集中度上升}}
\author{KC桌面 自动生成}
\date{260515}
\begin{document}
\maketitle
\section*{一页摘要}
\begin{itemize}
  \item 全球石油供应中断规模远超2022年，但通胀预期和信用利差异常平静，市场可能系统性低估了高油价的持续时间。
  \item 美股盈利增长极度依赖TECH+板块，剔除后增速从28.4\%骤降至11.1\%，结构性断层加深。
  \item AI算力供需失衡至少持续到2027年下半年，云厂商资本开支上调反映真实需求而非盲目投资。
  \item 欧洲天然气储气进度落后五年均值，若夏季注气持续偏慢，冬季前补库压力将推高价格。
  \item 中国信贷数据疲软背后是信用创造机制从银行信贷向债券市场切换，政策传导存在结构性堵点。
  \item 日本银行股定价逻辑正从利率预期转向贷款结构性增长，企业资本支出相关贷款强劲。
\end{itemize}
\section{宏观与利率}
\textbf{全球利率市场波动性主要来自央行反应函数的不确定性，而非通胀数据本身；供给冲击下各国央行政策路径分歧加大。}

\begin{itemize}
  \item 市场对油价冲击的定价集中在能源价格本身，但对能源成本向非能源部门成本结构、消费者信心的二阶传导缺乏重视，油价上行风险已被部分计入，但更广泛的供给侧通胀压力尚未被充分定价。
  \item 美联储内部鹰鸽分歧正从节奏之争演变为框架之争，市场对央行未来路径的定价仍停留在需求驱动型通胀的旧框架中。
  \item 英国国债收益率升至1998年以来最高，核心原因是低增长与高通胀的致命组合，以及供给结构和持有者结构的永久性重置，远期收益率中枢可能永久性上移。
  \item 欧洲央行鸽派声音增大，市场开始押注后置降息，但能源价格高企和薪资增长使政策路径不确定性上升。
  \item 日本银行股过去一个月跑输大盘，但基本面叙事正从利率预期驱动转向贷款结构驱动，企业资本支出相关贷款增长强劲，贷款量的结构性扩张成为盈利改善的新支撑。
\end{itemize}
\begin{figure}[htbp]
\centering
\includegraphics[width=0.88\linewidth]{figures/fig_007.jpg}
\caption{Exhibit 1 - +44(20)7051-1664 loic.mathys@gs.com GS International Exhibit 1: Gilt yields are the highest among major DM bond markets G4 30y yields, GS fitted yields <details> <summary>line</sum}
\end{figure}
\begin{figure}[htbp]
\centering
\includegraphics[width=0.88\linewidth]{figures/fig_008.jpg}
\caption{Exhibit 2 - | 2022 | 1.0 | 1.5 | 1.3 | 0.3 | | 2024 | 1.5 | 2.0 | 1.8 | 0.5 | | 2026 | 2.0 | 2.5 | 2.3 | 1.0 | </details> Source: GS Global Investment Research, GS FICC and Equities Exhibit 2:}
\end{figure}
{\small\color{refgray}\textbf{References:} [R001] 市场真正低估的不是油价，而是高油价持续的时间长度; [R005] 市场真正低估的不是英国通胀，而是英债供给结构的永久性改变; [R016] 市场真正低估的不是通胀，而是供给冲击下的央行反应函数重构; [R034] 日本银行股真正的博弈点不在利率，而在贷款结构性的质变}

\section{AI/算力与科技硬件}
\textbf{AI需求持续超预期，供给瓶颈成为未来三年核心定价变量，但盈利增长高度集中于少数巨头，板块内部分化加剧。}

\begin{itemize}
  \item AI计算资源供需失衡至少到2027年下半年才可能初步缓解，云厂商资本开支上调反映真实需求，市场对资本开支的恐惧可能被过度定价。
  \item AI光模块行业进入需求持续超预期、供给持续紧约束的双强周期，Lumentum CEO表示两个季度内可填满2028年全年产能，市场尚未充分计入供给短缺带来的结构性溢价。
  \item ABF基板收入即将超越2022年四季度历史峰值，但本轮驱动因素从消费电子补库存切换为AI服务器出货加速，可替代性大幅降低，ASP环比增长有望加速至7-10\%。
  \item 标普500盈利增长28.4\%中TECH+板块贡献17.3个百分点，剔除后增速骤降至11.1\%，六大科技巨头增速61.2\%，盈利增长极度集中。
  \item 散户正在系统性地从软件股撤离并加码半导体，软件板块拥挤度跌至历史低位而半导体持续攀升，这种分化可能触发风格切换。
\end{itemize}
\begin{figure}[htbp]
\centering
\includegraphics[width=0.88\linewidth]{figures/fig_044.jpg}
\caption{Exhibit 2 - We believe our forecasts are more bullish than consensus. Compared with LightCounting's 65\% growth forecast in 2026, we estimate shipment growth for 800G+1.6T to exceed 200\%. For 2}
\end{figure}
\begin{figure}[htbp]
\centering
\includegraphics[width=0.88\linewidth]{figures/fig_054.jpg}
\caption{Figure 2 - Expected Final Reported Growth </details> Source: Standard \& Poor's, Refinitiv, FactSet, UBS Expected 1Q EPS Growth is now 28.6\% YoY Figure 2: S\&P 500 EPS Growth - Impact from TECH}
\end{figure}
{\small\color{refgray}\textbf{References:} [R002] 市场真正低估的不是通胀，而是“繁荣循环”的自我强化与政治脆弱性; [R009] 市场真正低估的不是IPO抽血，而是4100亿美元被动买入的重新定价; [R018] ABF基板正在经历的不是周期复苏，而是供给侧的永久性再定价; [R025] 市场真正低估的不是AI需求，而是供给失衡的持续时长; [R031] 光模块市场真正被低估的不是需求，而是供给短缺的持续性与定价权转移; [R032] 市场正在软件和半导体之间押注一个不对称的再平衡; [R039] 标普500盈利的真正故事不在28.4\%的数字里，而在数字背后的结构性断层}

\section{能源与大宗商品}
\textbf{全球石油和天然气供给端收缩是结构性的而非暂时的，亚洲需求回归和炼化产能受损正在重塑定价体系。}

\begin{itemize}
  \item 全球约13\%的石油供应处于离线状态远超2022年的3\%，但通胀预期和信用利差异常平静，市场可能系统性低估了高油价持续的时间长度。
  \item 美国汽油市场急剧收紧，库存低于历史季节性中位数5\%，批发价比亚洲和欧洲贵15\%，零售价距历史最高仅差0.5美元/加仑，出口激增和炼厂结构调整是主因。
  \item 欧洲天然气储气进度落后五年均值12个百分点，若注入速率持续偏慢，冬季前填充率可能仅80\%，追赶式补库将推高价格。
  \item 亚洲LNG需求出现明确复苏信号，中国进口量从36百万吨/年反弹至48百万吨/年，若霍尔木兹海峡中断持续，3Q26 TTF价格或从44欧元升至65欧元/兆瓦时。
  \item 中东炼化产能受损规模超350万桶/日，是2025年俄罗斯无人机袭击高峰期的三倍以上，修复周期可能以年为单位，成品油出口净缺口约350万桶/日。
  \item 化肥市场利润分配结构正在重塑，氮肥Q2见顶，磷肥被成本压死，钾肥相对坚挺但空间有限，投资逻辑从看价格涨跌转向看成本曲线位置。
\end{itemize}
\begin{figure}[htbp]
\centering
\includegraphics[width=0.88\linewidth]{figures/fig_013.jpg}
\caption{Exhibit 4 - \textbackslash{}- The IEA supply beat relative to our balance estimate was driven primarily by Iran and the UAE, likely reflecting less binding storage constraints than expected due to untrackabl}
\end{figure}
\begin{figure}[htbp]
\centering
\includegraphics[width=0.88\linewidth]{figures/fig_021.jpg}
\caption{Figure 1 - As of 12 May, European gas storage is \$36\textbackslash{}\%\$ full, vs. \$43\textbackslash{}\%\$ a year ago and \$48\textbackslash{}\%\$ for the five-year average. Net injections over the past week amounted to 1.5bcm, below last year}
\end{figure}
{\small\color{refgray}\textbf{References:} [R001] 市场真正低估的不是油价，而是高油价持续的时间长度; [R007] 亚洲LNG需求正在回归，这才是下半年TTF真正的上行风险; [R008] 美国汽油市场正在成为全球原油定价的新锚点; [R015] 全球炼化产能正在经历一场未被充分定价的结构性断裂; [R017] 欧洲天然气市场的真正风险不是缺气，而是“补气”的时间窗口正在关闭; [R038] 化肥市场的真正拐点不在价格，而在利润分配结构的重新洗牌}

\section{中国信贷与地产}
\textbf{信贷收缩背后是信用创造机制从银行信贷向债券市场切换，政策传导存在结构性堵点，而非单纯需求不足。}

\begin{itemize}
  \item 4月人民币贷款存量罕见负增长，但企业债券融资净增4540亿元几乎翻倍，企业融资需求并未消失而是在向债券市场转移，信用创造渠道正在重新分配。
  \item 票据融资飙升至创纪录的1.24万亿元，本质是银行满足信贷考核指标的填充物，剔除后企业实际贷款需求可能比表面数据显示的更弱。
  \item 居民贷款减少7870亿元，按揭和消费信贷持续偏弱，居民仍在去杠杆，存款从银行流向理财产品。
  \item 宽松的流动性环境未有效转化为实体信用扩张，DR007已降至1.3\%低于政策利率，但低利率并未刺激出足够借贷需求，政策传导存在结构性阻塞。
  \item 香港长端利率定价权正从资金淤积转向供给与复苏，非金融债券供给回归和住宅地产复苏可能改变长端利率的供需方程。
\end{itemize}
\begin{figure}[htbp]
\centering
\includegraphics[width=0.88\linewidth]{figures/fig_005.jpg}
\caption{Exhibit 2 - 2. RMB loan stock (not seasonally adjusted) fell for the first time since July 2025, significantly below market expectations. The outstanding RMB loan growth slowed to \$5.6\textbackslash{}\%\$ yoy }
\end{figure}
\begin{figure}[htbp]
\centering
\includegraphics[width=0.88\linewidth]{figures/fig_017.jpg}
\caption{Figure 1 - Economist william-w.deng@UBS.com +852-2971 6765 Figure 1: Credit growth slowed, credit impulse turned more negative <details> <summary>line</summary> | Year | Credit impulse (LHS) }
\end{figure}
{\small\color{refgray}\textbf{References:} [R004] 市场真正低估的不是信贷需求，而是供给侧的被动收缩; [R010] 信用扩张的真正减速器不是需求，而是信用创造机制的切换; [R011] 香港长端利率的定价权正在从“资金淤积”转向“供给与复苏”; [R014] 四月信贷转负，市场真正低估的不是需求，而是政策传导的堵点}

\section{区域市场与地缘政治}
\textbf{中美元首会晤建立战略稳定框架为竞争划定底线，日本经济韧性来自内部再平衡，印度能源安全叙事转向成本可控。}

\begin{itemize}
  \item 中美元首会晤提出建设性、战略稳定的双边关系新框架，为未来三到五年双边互动设定底层操作系统，系统性脱钩的极端风险被降低。
  \item 日本2026年Q1 GDP上修至+1.3\%，但中东冲击的实质影响将在Q2显现，工资-消费-资本支出的自我支撑机制是经济韧性的核心。
  \item 印度正从颜色正确的能源叙事转向成本可控的能源叙事，煤制气成为战略选项，在考虑补贴后已具备成本优势。
  \item 人民币中间价定价机制正从被动反映市场波动转向主动管理预期锚，模型误差的持续偏倚指向政策意图的重塑。
  \item 日系零部件企业在华面临的不再是周期性销量问题，而是中国本土供应商在价格和开发速度上的结构性竞争力挑战。
\end{itemize}
\begin{figure}[htbp]
\centering
\includegraphics[width=0.88\linewidth]{figures/fig_009.jpg}
\caption{Exhibit 1 - Below, we summarize our outlook for the main demand components from Q2 onward. Private consumption: Consumer sentiment has deteriorated significantly since March, due to concerns o}
\end{figure}
\begin{figure}[htbp]
\centering
\includegraphics[width=0.88\linewidth]{figures/fig_048.jpg}
\caption{Exhibit 1 - The downtrend in domestic demand in China remains unchanged. While exports from China are increasing, there is no impact on Japanese companies at this stage. In terms of assumption}
\end{figure}
{\small\color{refgray}\textbf{References:} [R006] 日本经济真正的韧性不在GDP数字，而在工资与能源替代之间的再平衡; [R012] 市场真正低估的不是汇率点位，而是定价机制的结构性切换; [R013] 中美元首会晤的真正信号：战略稳定框架的建立比任何单点突破都更重要; [R021] 印度正在用煤炭重新定义能源安全，而市场还在用旧地图看印度; [R033] 日系零部件供应商的真正考验：不是销量下滑，而是中国本土对手正在重写竞争规则}

\section{风险偏好与资金流向}
\textbf{散户交易行为正在改变美股价格发现机制，资金从软件向半导体再平衡，但拥挤度分化已达历史极值。}

\begin{itemize}
  \item 散户持有约12万亿美元股票资产占市值10\%，但交易量占比达20\%，其杠杆使用和订单流结构正在重塑股票价格对基本面信息的反应效率。
  \item 散户正从软件股系统性撤离，微软从4月净买入第二变为5月净卖出第一，软件板块拥挤度跌至接近负值区间。
  \item 半导体板块持续吸金，DRAM相关股票买入量周环比翻4倍以上，拥挤度攀升至正值高位，与软件的剪刀差在过去十年仅出现数次。
  \item 杠杆融资违约率绝对水平回落，但违约形态正从传统破产转向困境交换，约50\%的困境交换会在3年内再次违约，回收率系统性走低。
  \item 三大AI公司IPO后锁定期满将催生约4100亿美元被动买入需求，自由流通股跳升更多是指数提供商的机械性重新分类而非内部人减持。
\end{itemize}
\begin{figure}[htbp]
\centering
\includegraphics[width=0.88\linewidth]{figures/fig_039.jpg}
\caption{Exhibit 25 - valuations and exhibit above-average volatility, even controlling for fundamentals like sales growth. Stocks with elevated retail trading activity have also exhibited greater under}
\end{figure}
\begin{figure}[htbp]
\centering
\includegraphics[width=0.88\linewidth]{figures/fig_046.jpg}
\caption{Figure 1 - 10/15 - Trading around Tariffs, AI and Earnings 10/8 - Back after Summer Break 10/1 - Back to Stocks, Thematic ETFs, GLD Rush 9/24 - Precious Metals, Cyclicals and Domestic 9/17 - }
\end{figure}
{\small\color{refgray}\textbf{References:} [R003] 杠杆融资违约的真正风险不在于违约率本身，而在于违约形态的结构性侵蚀; [R009] 市场真正低估的不是IPO抽血，而是4100亿美元被动买入的重新定价; [R028] 散户正在改变美国股市的定价机制，但机构投资者可能尚未完全理解这一变化; [R032] 市场正在软件和半导体之间押注一个不对称的再平衡}

\section{行业与产业趋势}
\textbf{AI正从替代就业叙事转向生产力杠杆，网络安全需求从合规驱动转向地缘政治驱动，数据中心投资地域分化加剧。}

\begin{itemize}
  \item 商业服务板块对AI替代就业的恐惧可能过度，历史数据显示技术革命从未导致长期失业，AI更多是生产力杠杆而非裁员按钮，蓝领派遣公司反而受益。
  \item OT网络安全需求正进入上行周期，Fortinet Q1营收超预期6.9\%中约3-4个百分点来自OT防火墙需求，地缘政治冲突和欧洲新监管框架是双重催化剂。
  \item 数据中心市场结构性分化加剧，核心都市区电力排队长达7年、建造成本每GW 12-15亿美元，准入门槛远高于乡村市场，头部运营商穿越长周期能力成为关键。
  \item 中国数据中心运营商正从出租机柜向出租算力迁移，GPUaaS毛利率高达80\%，商业模式跃迁的财务含义未被当前估值充分定价。
  \item 全球燃气轮机订单Q1创历史同期新高29GW，北美贡献近60\%，数据中心需求正从补充性负荷转变为主导性需求，拥有产能和定价权的企业将获得持续超额回报。
  \item 奢侈品市场正从追逐新客转向服务核心，趋势敏感型消费者退潮，核心高净值客户更关注稀有性和工艺，品牌必须学会收缩。
\end{itemize}
\begin{figure}[htbp]
\centering
\includegraphics[width=0.88\linewidth]{figures/fig_025.jpg}
\caption{EXHIBIT 1 - \# BUSINESS SERVICES SECTORS FACING THE AI ROLL-OUT \# IS THE STOCK MARKET OVERREACTING TO THE RISKS OF AI? Over the past year and a half, market concerns have been building about th}
\end{figure}
\begin{figure}[htbp]
\centering
\includegraphics[width=0.88\linewidth]{figures/fig_033.jpg}
\caption{EXHIBIT 1 - EXHIBIT 1: Comparing Tiers of Data Center Markets <table><tr><td></td><td>Major Metro</td><td>Minor Metro</td><td>Rural</td></tr><tr><td>Use Cases</td><td>Inference, Enterprise Col}
\end{figure}
{\small\color{refgray}\textbf{References:} [R019] 市场真正低估的不是AI替代就业，而是商业服务公司如何把AI转化为定价权; [R020] 航空技术股真正的风险不在中东，而在即将到来的运力削减; [R023] 网络安全需求正在从“合规驱动”转向“地缘政治驱动的供给侧重置”; [R024] 数据中心投资的地域分化：真正稀缺的不是容量，而是核心城市可建土地; [R027] 中国数据中心真正的增长引擎不是规模，而是从“出租机柜”到“出租算力”的商业模式迁移; [R030] 全球燃气轮机订单创纪录，但市场真正低估的是供给侧的结构性再定价; [R036] 奢侈品市场的真正拐点：不是需求萎缩，而是品牌必须学会“收缩”}

\section*{结语}
综合来看，当前市场主线围绕供给冲击的结构性持续性与AI需求超预期展开，但盈利增长和资金流向的集中度上升意味着分化加剧。投资者需关注央行反应函数重构、亚洲需求回归对能源定价的影响，以及中国信用创造机制切换带来的结构性变化。
\vfill
{\small\color{refgray}Personal reading notes and learning share only. Not investment advice.}
\end{document}
